\section{相关工作}

\subsection{拒答表征与方向干预}

拒答方向研究试图把安全对齐后的行为定位到可解释的隐藏子空间。Arditi 等人通过有害与无害提示的均值差寻找单一方向，并展示激活干预与秩一权重编辑~\cite{arditi2024refusal}。Heretic 延续这一机械干预路线，但研究重点从“是否存在方向”转向“如何自动选择方向位置、组件和层权重，并显式管理行为漂移”。本文不把自动化工程扩展描述成新的拒答表征定理。

多方向方法提醒我们不要把一维结果写成普遍规律。Piras 等人使用自组织映射从有害表示中提取多个原型，再相对无害中心构造多个拒答方向，并在其实验设置中优于单方向基线~\cite{piras2026som}。Heretic 固定提交没有实现该方法；两者的关系是可测试的后续扩展，而不是当前系统的隐藏能力。

\subsection{高效权重更新与自动搜索}

LoRA 通过冻结基础权重并注入低秩矩阵降低适配存储成本~\cite{hu2022lora}。Heretic 不训练这些矩阵，而是把解析得到的方向投影直接写入 A/B 因子；FULL 模式再用低秩分解近似行范数恢复后的完整差值。该用法保留 LoRA 的可切换与可合并特性，却需要单独讨论分解误差，不能直接继承微调论文中的性能结论。

Optuna 为动态搜索空间、采样与试验管理提供通用框架~\cite{akiba2019optuna}。Heretic 的贡献在其具体目标和搜索结构：按组件采样分段线性层核，同时最小化关键词拒答与首词元 KL，并向用户展示 Pareto 集。TPE 的统计性质来自优化框架，方向范围、归一化策略与 scorer 语义则完全由 Heretic 固定源码定义。
